% !TEX program = xelatex
% !TeX encoding = utf8
% !TeX spellcheck = pl-PL

%%%%%%%%%%%%%%%%%%%%%%%%%%%%%%%%%%%%%%%%%%%%%%%%%%%%%%%%%%%%%%%%%%%%%%%%%%%
% Wybierz rodzaj pracy dyplomowej oraz wydział
% Pick thesis type and faculty
%%%%%%%%%%%%%%%%%%%%%%%%%%%%%%%%%%%%%%%%%%%%%%%%%%%%%%%%%%%%%%%%%%%%%%%%%%%
\documentclass[thesis=inz,faculty=mel]{EE-dyplom} 

% thesis=[inz|mgr|bsc|msc]
%  * inz - praca inżynierska
%  * mgr - praca magisterska
%  * bsc - bachelor thesis
%  * msc - master thesis

% Skróty nazw wydziałów zgodne z domenami internetowymi
% Abbreviations of Faculties according to Internet subdomains
% faculty=[
%	arch,
%	gik,
%	ee,
%	wip
%	]

%%%%%%%%%%%%%%%%%%%%%%%%%%%%%%%%%%%%%%%%%%%%%%%%%%%%%%%%%%%%%%%%%%%%%%%%%%%
% Konfiguracja - do personalizacji
% Configuration - to be personalized
%%%%%%%%%%%%%%%%%%%%%%%%%%%%%%%%%%%%%%%%%%%%%%%%%%%%%%%%%%%%%%%%%%%%%%%%%%%
%\instytut{Instytut Badawczy Zjednoczonych Przedsiębiorstw Elektronowych}
\kierunek{Robotyka i Automatyka}
\specjalnosc{Robotyka}
\title{Sterowanie predykcyjne robotem równoległym na podstawie danych}
\engtitle{Data-enabled Predictive Control of a Parallel Robot}
\album{327709}
\author{Łukasz Pogorzelski}
\promotor{dr hab. inż. Paweł Malczyk, prof. uczelni}
\date{2026}
\longdate{2026-01-03}

%%%%%%%%%%%%%%%%%%%%%%%%%%%%%%%%%%%%%%%%%%%%%%%%%%%%%%%%%%%%%%%%%%%%%%%%%%%
% Streszczenie pracy i abstract.
% In case of thesis in English swap the order - English version goes first.
%%%%%%%%%%%%%%%%%%%%%%%%%%%%%%%%%%%%%%%%%%%%%%%%%%%%%%%%%%%%%%%%%%%%%%%%%%%
\streszczeniepracy{
Niniejsza praca dyplomowa porusza zagadnienie sterowania predykcyjnego opartego na danych (Data-enabled Predictive Control -- DeePC) w zastosowaniu do robota równoległego o dwóch stopniach swobody. Głównym celem pracy było zbadanie skuteczności metod sterowania, które eliminują konieczność jawnej identyfikacji modelu parametrycznego na rzecz wykorzystania surowych danych historycznych trajektorii wejścia-wyjścia.

W pierwszej części pracy opracowano model dynamiczny manipulatora typu pięciobok przegubowy, wykorzystując metodę przecięcia pętli kinematycznej (metoda cut-joint) i formułując równania we współrzędnych złączowych. Następnie poprawność modelu zweryfikowano z rozwiązaniem referencyjnym w specjalistycznym oprogramowaniu uzyskując wysoką zgodność rezultatów symulacji.

W części badawczej zaimplementowano i porównano dwa warianty algorytmu: metodę \mbox{Regularized DeePC}, wykorzystującą pełną historię danych oraz zmodyfikowaną metodę Selective DeePC, umożliwiającą lokalny wybór danych pomiarowych. Badania symulacyjne przeprowadzono w środowisku MATLAB dla zróżnicowanych trajektorii (lemniskata Gerona, wielobok), zarówno w warunkach idealnych, jak i w obecności szumu pomiarowego.

Analiza wyników wykazała fundamentalne różnice w działaniu obu metod. Algorytm \mbox{Regularized} DeePC cechuje się wysoką stabilnością i gładkością sygnału sterującego, jednak w warunkach zakłóceń wykazuje tendencję do generowania uchybów statycznych i ,,ścinania'' ostrych zakrętów. Z kolei metoda Selective DeePC zapewnia lepsze odwzorowanie geometrii trajektorii i szybszą reakcję układu, co odbywa się jednak kosztem mniejszej stabilności i występowania oscylacji w sygnale sterującym.

Kluczowym wnioskiem jest wykazanie przewagi metody selektywnej pod względem wydajności obliczeniowej. Algorytm Selective DeePC pozwolił na siedmiokrotne skrócenie czasu obliczeń w stosunku do metody klasycznej (redukcja ze 100 ms do ok. 14 ms), co czyni go perspektywicznym rozwiązaniem dla systemów czasu rzeczywistego.

}

\slowakluczowe{robot równoległy, współrzędne złączowe, regulacja predykcyjna, DeePC, \mbox{Data-enabled} Predictive Control}

\thesisabstract{
This thesis addresses the issue of Data-enabled Predictive Control (DeePC) applied to a parallel robot with two degrees of freedom. The main objective of the thesis was to investigate the effectiveness of control methods that eliminate the need for explicit parametric model identification in favor of using raw historical input-output data trajectories.

In the first part of the thesis, a dynamic model of a five-bar linkage manipulator was developed using the cut-joint method in joint coordinates. Subsequently, the model's correctness was validated against a reference solution in specialized software, achieving high simulation consistency.

In the research part, two variants of the algorithm were implemented and compared: the classic Regularized DeePC, utilizing the full data history with regularization, and the modified Selective DeePC, based on a local data selection mechanism. Simulation studies were conducted in the MATLAB environment for various trajectories (Gerono lemniscate, polygon), both under ideal conditions and in the presence of measurement noise.

Analysis of the results revealed fundamental differences in the performance of both methods. The Regularized DeePC algorithm is characterized by high stability and smoothness of the control signal; however, under disturbance conditions, it tends to generate static errors and "cut" sharp corners. In turn, the Selective DeePC method ensures better trajectory geometry tracking and faster system response, but at the cost of lower stability and the occurrence of oscillations in the control signal.

A key conclusion is the demonstration of the selective method's advantage in terms of computational efficiency. The Selective DeePC algorithm allowed for a sevenfold reduction in computation time compared to the classic method (a reduction from 100 ms to approx. 14 ms), making it a promising solution for real-time systems.} % end of abstract

\thesiskeywords{parallel robot, joint coordinates, predictive control, DeePC, Data-enabled Predictive Control}
%%%%%%%%%%%%%%%%%%%%%%%%%%%%%%%%%%%%%%%%%%%%%%%%%%%%%%%%%%%%%%%%%%%%%%%%%%%
% Tu zaczyna się dokument
% Here is the beginning of the document
%%%%%%%%%%%%%%%%%%%%%%%%%%%%%%%%%%%%%%%%%%%%%%%%%%%%%%%%%%%%%%%%%%%%%%%%%%%
\begin{document}
    % Strony nagłówkowe - musi być
    % Headers - must exist
    %\frontpages
    
    % Tekst pracy - osobne pliki do edycji
    % Thesis text - separate files for edition

	\chapter{Introduction}
\label{ch:wstep}

\section{Introduction and Background}

Modern robotics imposes increasingly stringent requirements on control systems regarding precision and adaptation to uncertain operating conditions. For years, Model Predictive Control (\textbf{MPC}) has been the industrial standard. Despite its effectiveness, classical MPC possesses a significant drawback: it requires an accurate mathematical model of the controlled plant. In the case of complex nonlinear systems, such as parallel robots, the identification of physical parameters is a difficult, time-consuming process, often prone to errors that negatively affect control performance.

These limitations have motivated the development of \textit{Data-Driven Control} methods. This paradigm assumes bypassing the explicit stage of modeling the object's physics in favor of utilizing raw historical data (input-output trajectories). One of the most promising algorithms in this field is \textbf{DeePC} (\textit{Data-Enabled Predictive Control}), introduced in 2019 \cite{deepc}. This method combines the mathematical rigor of classical prediction with the flexibility of a model-free approach.

However, a key challenge in the DeePC method is the appropriate selection of data and the ability to control nonlinear systems in the presence of noise occurring in measurement channels. The literature proposes various strategies to address these issues—ranging from methods based on strong regularization (\textbf{Regularized DeePC}) to methods based on local data selection (\textbf{Selective DeePC}, also known in the literature as \textit{Select-DeePC} \cite{choosewisely}). This thesis focuses on investigating and comparing these two strategies in the context of controlling a robot with complex kinematics, as representative approaches fitting within the current state of knowledge regarding data-driven methods.

\section{Positioning the DeePC Method Against Model-Based and Data-Driven Solutions}

Choosing an appropriate control strategy for robotic manipulators involves balancing the required knowledge of the system (model) and the amount of data necessary to tune the controller. This problem is typically considered in the literature within two extreme paradigms: model-based control and data-driven control \cite{deepc, tatiewski, quad, ethzurich}.

\subsection{Limitations of Model-Based Methods (MPC)}
One of the dominant solutions in advanced control systems is classical MPC, which allows for the explicit inclusion of constraints on control signals and state variables, thereby ensuring a high level of operational safety \cite{nmpc}. However, the effectiveness of MPC is highly dependent on the quality of the available parametric model of the plant \cite{ethzurich}.

These limitations constitute a significant motivation for developing methods that retain the advantages of predictive control while eliminating the need for explicit modeling of the plant physics. This leads toward data-driven approaches, such as the DeePC algorithm \cite{deepc}.

\subsection{Challenges of Pure Data-Driven Methods (RL)}
At the opposite end of the spectrum are Reinforcement Learning (RL) methods, which completely forgo a parametric model in favor of optimizing a reward-based control policy. Although this approach eliminates the system identification problem, several critical drawbacks arise in robotic applications \cite{deepc, ethzurich}:
\begin{itemize}
    \item \textbf{Lack of safety guarantees:} RL methods inherently do not account for hard state and control constraints in a deterministic manner, which is crucial in MPC.
    \item \textbf{Low sample efficiency:} These algorithms require massive training datasets, which is time-consuming and costly in real-world experiments.
    \item \textbf{Lack of reproducibility:} Results are often highly sensitive to the choice of hyperparameters and the random exploration process, making stability verification difficult.
\end{itemize}

\subsection{DeePC as a Synthesis of MPC and Data-Driven Advantages}
The \textit{Data-enabled Predictive Control} algorithm was chosen as the subject of research in this thesis because it bridges the gap between the aforementioned approaches. Unlike RL, DeePC is based on a receding horizon predictive control structure, allowing for the explicit inclusion of safety constraints (similar to MPC) \cite{nmpc, deepc, ethzurich}.

The key innovation of DeePC, however, is the replacement of the parametric model (required in methods described, e.g., in \cite{tatiewski, nmpc}) with Hankel matrices constructed from raw input-output trajectories. Based on Willems' Fundamental Lemma and behavioral systems theory, DeePC allows for predicting the future behavior of the system directly from historical data \cite{deepc, WILLEMS2005325}.

\section{Aim and Scope of the Thesis}

The main objective of this diploma thesis is the implementation and comparative analysis of two variants of the data-driven predictive control algorithm (\textbf{DeePC}) applied to the position control of a parallel robot with two degrees of freedom.

The following were analyzed:
\begin{itemize}
    \item \textbf{Regularized DeePC:} The classical approach, utilizing the full available data history with the application of regularization terms for noise filtration.
    \item \textbf{Selective DeePC:} A modified approach, utilizing a data selection mechanism (local, implicit linearization) to more accurately represent the object's behavior at the current operating point.
\end{itemize}

\textbf{The specific objectives of the thesis include:}
\begin{enumerate}
    \item Development of a simulation model of a parallel robot (five-bar linkage type) in joint coordinates and its validation.
    \item Implementation of the \textbf{Regularized DeePC} algorithm and the \textbf{Selective DeePC} algorithm in the MATLAB environment.
    \item Conducting simulation studies for both methods under ideal conditions (no noise) and in the presence of measurement disturbances, for various motion trajectories (lemniscate, polygon).
    \item Sensitivity analysis of the algorithms to key hyperparameters (prediction/initialization horizons, regularization coefficients, number of samples).
    \item Comparison of both methods in terms of tracking precision, stability, control signal characteristics, and computational complexity.
\end{enumerate}

\section{System Architecture}
\label{sec:architektura_systemu}

To facilitate the analysis of the subsequent parts of the thesis and to understand the interactions occurring between individual elements of the project, a conceptual diagram presented in Figure \ref{fig:najwazniejszy_graf} was developed. It illustrates the signal flow in a closed-loop feedback system with the following components:
\begin{itemize}
    \item \textbf{Reference trajectory} defines the expected behavior of the robot ($\boldsymbol{y}_{ref}$), serving as the target for the control system.
    \item \textbf{Predictive controller (DeePC)} performs the decision-making role. Using the historical data library (represented by the matrix $\mathcal{H}$) and current feedback measurements, it determines the optimal control signal ($\boldsymbol{u}_{opt}$).
    \item \textbf{Dynamic model} constitutes the virtual control plant. It receives the calculated torques and then—based on the equations of motion—determines the mechanical response of the system (accelerations $\boldsymbol{\ddot \theta}$).
    \item \textbf{Simulation environment (Integrator)} transforms the system of equations into motion, updating the robot's state (positions and velocities—$\boldsymbol{\theta}$ and $\boldsymbol{\dot \theta}$) and generating a new measurement vector ($\boldsymbol{y}$), which closes the control loop.
\end{itemize}

\begin{figure}[!htbp]
    \centering
    \includegraphics[width=1\linewidth]{rysunki/1.wstep/Graf_cala_praca (3) (3).drawio.pdf}
    \caption{Conceptual diagram of the implemented control system architecture}
    \label{fig:najwazniejszy_graf}
\end{figure}
    \chapter{Dynamics Analysis of the Parallel Robot}
\label{ch:model_dynamiczny}

This chapter presents the mathematical model that serves as the foundation for the synthesis of predictive controllers. The subject of research is a parallel robot structure with design parameters of the Quanser "2 DOF ROBOT" \cite{quanser, kkr}, which is a planar five-bar linkage mechanism characterized by nonlinear dynamics. The robot dynamics are described in joint coordinates using the cut-joint method \cite{nikravesh}.

\section{General Formulation of the Dynamics Problem}
\label{sec:teoria_ogolna}

Following \cite{nikravesh}, the dynamics of the closed-loop mechanism are formulated by virtually cutting the kinematic loop, yielding a set of open chains governed by holonomic constraints:
\begin{equation}
    \label{eq:general_loop_constraint}
    \boldsymbol{\Phi}(\boldsymbol{q}) = \boldsymbol{0}
\end{equation}
Differentiating with respect to time yields the constraint at the velocity level:
\begin{equation}
    \label{eq:D_matrix_beg}
    \mathbf{D}\boldsymbol{\dot q} = \boldsymbol{0}
\end{equation}
where $\mathbf{D} = \frac{\partial \boldsymbol{\Phi}}{\partial \boldsymbol{q}}$. Introducing the velocity transformation matrix $\mathbf{B}$ such that $\boldsymbol{\dot q}= \mathbf{B}\boldsymbol{\dot \theta}$, the joint-space Jacobian is defined as:
\begin{equation}
    \label{eq:C_matrix_beg}
    \mathbf{C}=\mathbf{D}\mathbf{B}
\end{equation}
and the acceleration-level constraint follows as:
\begin{equation}
    \label{eq:acceleration_constraint}
    \mathbf{C}\boldsymbol{\ddot \theta} + \mathbf{\dot C}\boldsymbol{\dot \theta}= \boldsymbol{0}
\end{equation}
The equations of motion with constraint reaction forces (Lagrange multipliers $\boldsymbol{\prescript{\star}{}{}\lambda}$) take the DAE form:
\begin{equation}
    \label{eq:general_DAE}
    \begin{bmatrix}
        \boldsymbol{M}_{joint}&-\mathbf{C}^\mathrm{T}\\
        \mathbf{C}&\boldsymbol{0}
    \end{bmatrix}
    \begin{bmatrix}
        \boldsymbol{\ddot\theta}\\
        \prescript{\star}{}{}\boldsymbol{\lambda}
    \end{bmatrix}
    =
    \begin{bmatrix}
        \prescript{(a)}{}{\boldsymbol{h}_{joint}}\\
        -\mathbf{\dot C}\boldsymbol{\dot \theta}
    \end{bmatrix}
\end{equation}
where $\boldsymbol{M}_{joint} = \mathbf{B}^\mathrm{T}\boldsymbol{M}\mathbf{B}$ is the mass matrix projected onto the joint space, and $\prescript{(a)}{}{\boldsymbol{h}_{joint}}$ is the vector of generalized forces in the joint space.

\section{Detailed Model of the Parallel Robot}
\label{subch:implementacja_robot}

The geometry of the mechanism with the marked virtual cut point (point C), adopted coordinate systems, and notation follow \cite{nikravesh}; key vectors are illustrated in Figure \ref{fig:wybrane_wektory}. The kinematic scheme with joint variables and the operating trajectory is shown in Figure \ref{fig:robot_rownolegly}.

\begin{figure}[!htbp]
    \centering
    \includegraphics[width=1\linewidth]{rysunki/2.model_dynamiczny/wybrane_wektory.pdf}
    \caption{Representation of selected vectors on the left half of the mechanism (formed by points ABC) of the parallel robot}
    \label{fig:wybrane_wektory}
\end{figure}

\begin{figure}[!htbp]
    \centering
    \includegraphics[width=1\linewidth]{rysunki/2.model_dynamiczny/robot_rownolegly.pdf}
    \caption{Kinematic scheme of the robot including coordinate systems, joint variables, and the operating trajectory (Trajectory A)}
    \label{fig:robot_rownolegly}
\end{figure}

Direct torque control is assumed, neglecting the dynamics of the drive systems, which is a sufficient simplification for the synthesis of predictive control algorithms.

\subsection{Definition of Coordinates and Constraints}
The state variables of the system were defined as follows:
\begin{itemize}
    \item \textbf{Absolute coordinate vector} $\boldsymbol{q}$ (12 variables) -- contains the positions of the centers of mass $\boldsymbol{r}_i$ and orientation angles $\varphi_i$ for each of the four links:
    \begin{equation}
        \label{eq:vector_c}
        \boldsymbol{q} =\begin{bmatrix}
            \boldsymbol{r}_1^\mathrm{T} & \varphi_1 & \boldsymbol{r}_2^\mathrm{T} & \varphi_2 & \boldsymbol{r}_3^\mathrm{T} & \varphi_3 & \boldsymbol{r}_4^\mathrm{T} & \varphi_4
        \end{bmatrix}^\mathrm{T}
    \end{equation}
    \item \textbf{Joint coordinate vector} $\boldsymbol{\theta}$ (4 variables) -- describes the configuration of the two open kinematic trees formed after the virtual cut at point C:
    \begin{equation}
        \label{eq:vector_theta}
        \boldsymbol{\theta} = \begin{bmatrix}
            \theta_1 & \theta_2 & \theta_3 & \theta_4
        \end{bmatrix}^\mathrm{T}
    \end{equation}
\end{itemize}

The constraint equation $\prescript{\star}{}{}\boldsymbol{\Phi}$ results from the loop closure condition at point C:
\begin{equation}
    \label{eq:Phi}
    \begin{aligned}
        \prescript{\star}{}{}\boldsymbol{\Phi}(\boldsymbol{q}) = \underbrace{\boldsymbol{r}_2 + \boldsymbol{s}_2^C}_{\text{ABC Chain}} &- \underbrace{(\boldsymbol{r}_3 + \boldsymbol{s}_3^C)}_{\text{EDC Chain}} = \boldsymbol{0} \\
        \prescript{\star}{}{}\boldsymbol{\Phi}(\boldsymbol{q}) = (-\boldsymbol{s}_1^A + \boldsymbol{s}_1^B - \boldsymbol{s}_2^B + \boldsymbol{s}_2^C) &- (-\boldsymbol{s}_3^D + \boldsymbol{s}_4^D - \boldsymbol{s}_4^E + \boldsymbol{L}) - \boldsymbol{s}_3^C = \boldsymbol{0}
    \end{aligned}
\end{equation}
where $\boldsymbol{L} = [2l, 0]^\mathrm{T}$ is the distance vector between the motor axes.

\subsection{Determination of Structural Matrices}
Following the general procedure of Section \ref{sec:teoria_ogolna}, the matrices necessary to construct the equations of motion (\ref{eq:general_DAE}) were determined.

\textbf{Velocity transformation matrix ($\mathbf{B}$)}:
By analyzing the topology of the open chains (revolute joints $\boldsymbol{R}_{ij}$), the matrix $\mathbf{B}$ was constructed:
\begin{equation}
    \label{eq:B_matrix}
    \mathbf{B} = \begin{bmatrix}
        \boldsymbol{R_{11}} & \boldsymbol{0}_{3\times1} &\boldsymbol{0}_{3\times1} &\boldsymbol{0}_{3\times1}\\
        \boldsymbol{R}_{21} & \boldsymbol{R}_{22} & \boldsymbol{0}_{3\times1} & \boldsymbol{0}_{3\times1}\\ 
        \boldsymbol{0}_{3\times1} & \boldsymbol{0}_{3\times1}&\boldsymbol{R}_{33}&\boldsymbol{R}_{34}\\
        \boldsymbol{0}_{3\times1} & \boldsymbol{0}_{3\times1} & \boldsymbol{0}_{3\times1} & \boldsymbol{R}_{44}
     \end{bmatrix}_{12\times4}
\end{equation}
The elements $\boldsymbol{R}_{ij}$ depend on the vectors $\boldsymbol{\breve d}_{ij}$; their detailed definitions and the form of $\mathbf{\dot B}$ are included in \textbf{Appendix \ref{app:szczegoly_kinematyczne}}.

\textbf{Constraint Jacobian matrix ($\mathbf{D}$)}:
Differentiating the loop equation (\ref{eq:Phi}) yields:
\begin{equation}
    \label{eq:vector_loop_differentiated}
    \prescript{\star}{}{\boldsymbol{\dot \Phi}} = \boldsymbol{\dot r}_2 + \boldsymbol{\breve s}_2^C\dot \varphi_2-\boldsymbol{\dot r}_3- \boldsymbol{\breve s}_3^C\dot \varphi_3
\end{equation}
from which the explicit form of $\prescript{\star}{}{\mathbf{D}}$ follows:
\begin{equation}
    \label{eq:absolute_jacobain}
        \prescript{\star}{}{\mathbf{D}}
        = 
        \begin{bmatrix}
        \boldsymbol{0}_{2\times2} &  \boldsymbol{0}_{2\times2} & \boldsymbol{I}_{2\times2} &  \boldsymbol{\breve s}_2^C & \boldsymbol{-I}_{2\times2} & -\boldsymbol{\breve s}_3^C &  \boldsymbol{0}_{2\times2} &  \boldsymbol{0}_{2\times2}
        \end{bmatrix}
\end{equation}

\textbf{Joint Space Jacobian matrix ($\mathbf{C}$)}:
\begin{equation}
    \label{eq:c_dot_matrix}
    \mathbf{C}=\prescript{\star}{}{\mathbf{D}}\mathbf{B}, \quad \quad \mathbf{\dot C}= \prescript{\star}{}{\mathbf{D}}\mathbf{\dot B} + \prescript{\star}{}{\mathbf{\dot D}}\mathbf{B}
\end{equation}

\subsection{Final System of Dynamics Equations}
Adopting a diagonal mass matrix $\boldsymbol{M}$ (according to Table \ref{tab:parametry_czlonu}) and a vector of applied external forces $\prescript{(a)}{}{\boldsymbol{h}}$ containing control torques $\tau_1, \tau_4$:
\begin{equation}
    \label{eq:mass_matrix}
        \boldsymbol{M} = \boldsymbol{diag}(m, m,J_z, \dots, m, m, J_z)
\end{equation}
\begin{equation}
    \label{eq:h_applied}
        \prescript{(a)}{}{\boldsymbol{h}} = \begin{bmatrix}
            0 & 0 & \tau_1 &\boldsymbol{0}_{1\times3} &\boldsymbol{0}_{1\times3}&0&0 &\tau_4
        \end{bmatrix}^\mathrm{T}
\end{equation}
the joint-space projections are determined as:
\begin{equation}
    \label{eq:mass_matrix_joint}
        \boldsymbol{M}_{joint} = \mathbf{B}^\mathrm{T}\boldsymbol{M}\mathbf{B}
\end{equation}
\begin{equation}
    \label{eq:h_applied_joint}
        \prescript{(a)}{}{\boldsymbol{h}}_{joint} = \mathbf{B}^\mathrm{T}(\prescript{(a)}{}{\boldsymbol{h}} -\boldsymbol{M}\mathbf{\dot B}\boldsymbol{\dot \theta})
\end{equation}

To counteract constraint drift, Baumgarte stabilization \cite{BAUMGARTE19721} was applied by augmenting the acceleration constraint (\ref{eq:acceleration_constraint}):
\begin{equation}
    \label{eq:baumgarte}
    -\mathbf{\dot C}\boldsymbol{\dot \theta}_{new} = -\mathbf{\dot C}\boldsymbol{\dot \theta} -2\alpha \prescript{\star}{}{\boldsymbol{\dot\Phi}}
    -\beta^2\prescript{\star}{}{\boldsymbol{\Phi}}
\end{equation}
Parameters $\alpha = \beta = 10$ were selected heuristically and effectively suppress drift to $\sim10^{-11}$ m. The model was validated against MSC Adams simulations, showing high agreement; residual differences are attributed to differing numerical integration schemes.

\section{Model Parameters}
The geometric dimensions and inertial parameters adopted for the simulation are collected in Table \ref{tab:parametry_czlonu} (based on \cite{quanser, kkr}).

\begin{table}[!htbp]
    \centering
    \caption{Physical parameters of a single link}
    \label{tab:parametry_czlonu}
    \begin{tabular}{ccc}\toprule
        \textbf{Quantity} & \textbf{Symbol} & \textbf{Value} \\\midrule
        Link length & $l$ & $0.127\,\mathrm{m}$\\
        Link mass & $m$ & $0.065\,\mathrm{kg}$ \\
        Moment of inertia & $J_z$ & $9 \cdot 10^{-5}\,\mathrm{kg\cdot m^2}$ \\\bottomrule
    \end{tabular}
\end{table}
    \chapter{Synthesis of the Predictive Controller}
\label{ch:deepc}

\section{Classical Model Predictive Control -- MPC}
\label{subs:mpc}

Classical MPC minimizes a receding-horizon cost function subject to system dynamics and actuator constraints \cite{nmpc}. For a linear time-invariant (LTI) system described by:
\begin{align}
    \boldsymbol{x}_{k+1} &= \boldsymbol{A}\boldsymbol{x}_k + \boldsymbol{B} \boldsymbol{u}_k 
    \label{eq:statempc}
    \\
    \boldsymbol{y}_k &= \boldsymbol{C} \boldsymbol{x}_k + \boldsymbol{D} \boldsymbol{u}_k
    \label{eq:outmpc}
\end{align}
where $\boldsymbol{x} \in \mathbb{R}^n$, $\boldsymbol{u} \in \mathbb{R}^m$, $\boldsymbol{y} \in \mathbb{R}^p$, the optimization problem over prediction horizon $N$ takes the form:
\begin{equation}
    \min_{\boldsymbol{u},\boldsymbol{x},\boldsymbol{y}} \sum_{k=0}^{N-1} \left( \left\| \boldsymbol{y}_k - \boldsymbol{r}_{k} \right\|_{\boldsymbol{Q}}^2 + \left\| \boldsymbol{u}_k \right\|_{\boldsymbol{R}}^2 \right)
    \label{eq:costmpc}
\end{equation}
subject to:
\begin{align}
    \boldsymbol{u}_k &\in \mathcal{U}, \quad \forall k \in \{0, \ldots, N-1\}
    \label{eq:umpc}
    \\
    \boldsymbol{y}_k &\in \mathcal{Y}, \quad \forall k \in \{0, \ldots, N-1\}
    \label{eq:ympc}
\end{align}

Despite its effectiveness, MPC requires a parametric model identified in a separate offline stage — a process that does not guarantee optimality from the control perspective \cite{markovsky2023data}. The DeePC algorithm eliminates this stage by operating directly on raw measurement data as a non-parametric behavioral representation of the system.

\section{Foundations of Behavioral Control Theory}

The foundation of the DeePC algorithm is the behavioral approach to systems theory, in which a dynamic system is defined as the set of all permissible input-output trajectories \cite{WILLEMS2005325}. Key concepts enabling such a representation are defined below \cite{quad, ethzurich, markovsky2023data, deepc}.

\textbf{Hankel Matrix} $\mathcal{H}$ -- for a discrete signal $\boldsymbol{w} = \begin{bmatrix} \boldsymbol{w}_1 & \boldsymbol{w}_2 & \dots & \boldsymbol{w}_T \end{bmatrix}^\mathrm{T}$ of length $T$, the Hankel matrix of depth $L = T_{ini} + N$ is defined as:
\begin{equation}
    \label{eq:hankel_matrix}
    \mathcal{H}_L(\boldsymbol{w}) = 
    \begin{bmatrix}
        \boldsymbol{w}_1 & \boldsymbol{w}_2 & \cdots & \boldsymbol{w}_{T-L+1}\\
        \boldsymbol{w}_2 & \boldsymbol{w}_3 & \cdots & \boldsymbol{w}_{T-L+2}\\
        \vdots & \vdots & \ddots & \vdots\\
        \boldsymbol{w}_L & \boldsymbol{w}_{L+1} & \cdots & \boldsymbol{w}_T
    \end{bmatrix}
\end{equation}

\textbf{Mosaic Hankel Matrix} -- for $K$ independent datasets $\boldsymbol{w}^{d_{1}}, \dots, \boldsymbol{w}^{d_{K}}$, the mosaic matrix is formed by horizontal concatenation:
\begin{equation}
    \label{eq:mosaic_hankel}
    \mathcal{H}_{mos}(\boldsymbol{w}^d) = 
    \begin{bmatrix}
        \mathcal{H}_L(\boldsymbol{w}^{d_{1}}) & \mathcal{H}_L(\boldsymbol{w}^{d_{2}}) & \cdots & \mathcal{H}_L(\boldsymbol{w}^{d_{K}})
    \end{bmatrix}
\end{equation}

\textbf{Persistency of Excitation (PE)} -- an input signal $\boldsymbol{u}$ is sufficiently exciting of order $L$ if $\mathcal{H}_L(\boldsymbol{u})$ has full row rank \cite{ethzurich}. For a system of order $n$, the input must be sufficiently exciting of order $T_{ini} + N + n$, requiring the total dataset length to satisfy:
\begin{equation}
    T \ge (m+1)(T_{ini} + N + n) - 1
\end{equation}

\textbf{Willems' Fundamental Lemma} -- for an LTI system with a sufficiently exciting input of order $T_{ini} + N + n$, the columns of $\mathcal{H}_L$ span the space of all possible input-output trajectories of length $L = T_{ini} + N$ \cite{deepc, WILLEMS2005325}. Consequently, any future system trajectory can be expressed as a linear combination of the columns of this matrix, without identifying the state matrices $\boldsymbol{A}, \boldsymbol{B}, \boldsymbol{C}, \boldsymbol{D}$.

\section{Data-Driven Control Algorithm -- DeePC}
\label{sec:deepc_algo}

The non-parametric model of the system is constructed from historical data $\boldsymbol{u}^d$, $\boldsymbol{y}^d$ of length $T$. The Hankel matrix is partitioned along the row dimension into past (index $p$, depth $T_{ini}$) and future (index $f$, depth $N$) blocks:
\begin{equation}
    \label{eq:data_matrix_deepc}
    \mathcal{H} = \begin{bmatrix}
        \mathcal{H}_{p} \\ \hline \mathcal{H}_{f}
    \end{bmatrix} = 
    \begin{bmatrix}
         \begin{pmatrix} \boldsymbol{U}_p \\ \boldsymbol{Y}_p \end{pmatrix} \\[1em]
         \hline
         \begin{pmatrix} \boldsymbol{U}_f \\ \boldsymbol{Y}_f \end{pmatrix}
    \end{bmatrix}
\end{equation}
as illustrated in Figure \ref{fig:deepc_matrix_structure}.

\begin{figure}[!htbp]
    \centering
    \includegraphics[width=0.6\linewidth]{rysunki/3.deepc/graf_deepc (3).drawio.pdf}
    \caption{Data matrix structure used in the DeePC algorithm. The gray area includes $T_{ini}$ rows responsible for establishing the initial state (past), while the white area contains $N$ rows defining possible future evolutions of the system. Each column represents one physically realizable system trajectory of length $L = T_{ini} + N$.}
    \label{fig:deepc_matrix_structure}
\end{figure}

\subsection{Trajectory Prediction and the Decision Variable $\boldsymbol{g}$}

A decision vector $\boldsymbol{g} \in \mathbb{R}^{T-L+1}$ is introduced, weighting the columns of the data matrix. At each control iteration, the sought trajectory must satisfy:
\begin{equation}
\label{eq:deePC_regression}
\begin{bmatrix}
    \boldsymbol{U}_p \\
    \boldsymbol{Y}_p \\
    \boldsymbol{U}_f \\
    \boldsymbol{Y}_f
    \end{bmatrix}
    \boldsymbol{g} =
    \begin{bmatrix}
    \boldsymbol{u}_{ini} \\
    \boldsymbol{y}_{ini} \\
    \boldsymbol{u} \\
    \boldsymbol{y}
    \end{bmatrix}
\end{equation}
where $\boldsymbol{u}_{ini}$, $\boldsymbol{y}_{ini}$ contain the last $T_{ini}$ measured samples (with $T_{ini} \ge n$ to uniquely define the initial state \cite{ethzurich}), and $\boldsymbol{u}$, $\boldsymbol{y}$ are the future control and output sequences over horizon $N$.

\subsection{Formulation of the Optimization Problem}

The DeePC control problem is formulated as a quadratic programming (QP) task:
\begin{equation}
    \label{eq:deepc_optimization}
    \begin{aligned}
    \min_{\boldsymbol{g}, \boldsymbol{u}, \boldsymbol{y}} \quad & \sum_{k=0}^{N-1} \left( \| \boldsymbol{y}_k - \boldsymbol{r}_{k} \|_{\boldsymbol{Q}}^2 + \| \boldsymbol{u}_k \|_{\boldsymbol{R}}^2 \right) \\
    \text{s.t.} \quad & \begin{bmatrix} \boldsymbol{U}_p \\ \boldsymbol{Y}_p \\ \boldsymbol{U}_f \\ \boldsymbol{Y}_f \end{bmatrix} \boldsymbol{g} = \begin{bmatrix} \boldsymbol{u}_{ini} \\ \boldsymbol{y}_{ini} \\ \boldsymbol{u} \\ \boldsymbol{y} \end{bmatrix} \\
    & \boldsymbol{u}_k \in \mathcal{U}, \quad \boldsymbol{y}_k \in \mathcal{Y}
    \end{aligned}
\end{equation}
For LTI systems without noise, this formulation is mathematically equivalent to classical MPC \cite{ethzurich}. The algorithm operates in receding horizon mode: at each step $t$, the optimization (\ref{eq:deepc_optimization}) is solved and only the first element $\boldsymbol{u}_{opt}(t) = \boldsymbol{u}^*_0 = \boldsymbol{U}_f\boldsymbol{g}^*$ is applied to the plant (Figure \ref{fig:algorytm_deepc}).

\begin{figure}[!htbp]
    \centering
    \includegraphics[width=1\linewidth]{rysunki/3.deepc/algorytm.drawio (2).pdf}
    \caption{Block diagram of the DeePC controller operation in a closed loop.}
    \label{fig:algorytm_deepc}
\end{figure}

\section{Robust Control for Noise and Nonlinearities -- Regularized DeePC}

To address measurement noise and plant nonlinearities, the \textbf{Regularized DeePC} variant \cite{deepc, ethzurich} introduces a slack variable $\boldsymbol{\sigma}$ and regularization terms into the cost function:
\begin{equation}
    \label{eq:regularized}
    \min_{\boldsymbol{g}, \boldsymbol{u}, \boldsymbol{y}, \boldsymbol{\sigma}} \sum_{k=0}^{N-1} \left( \| \boldsymbol{y}_k - \boldsymbol{r}_{k} \|_{\boldsymbol{Q}}^2 + \| \boldsymbol{u}_k \|_{\boldsymbol{R}}^2 \right) + \lambda_g \|\boldsymbol{g}\|_2^2 + \lambda_{\sigma} \|\boldsymbol{\sigma}\|^2_2
\end{equation}
where $\lambda_g$ prevents overfitting to measurement noise and improves numerical conditioning \cite{ethzurich, markovsky2023data}, and $\lambda_{\sigma}$ controls the degree of trust in current measurements. The slack variable modifies the initialization constraint:
\begin{equation}
\begin{bmatrix}
    \boldsymbol{U}_p \\
    \boldsymbol{Y}_p \\
    \boldsymbol{U}_f \\
    \boldsymbol{Y}_f
    \end{bmatrix}
    \boldsymbol{g} =
    \begin{bmatrix}
    \boldsymbol{u}_{ini} \\
    \boldsymbol{y}_{ini} + \boldsymbol{\sigma} \\
    \boldsymbol{u} \\
    \boldsymbol{y}
    \end{bmatrix}
    \label{eq:regularized_set}
\end{equation}
allowing the linear combination of historical data to reproduce the measurement $\boldsymbol{y}_{ini}$ within the accuracy of $\boldsymbol{\sigma}$. Quadratic ($l_2$) regularization was adopted, keeping the problem as a standard QP task while ensuring practical exponential stability in the presence of noise \cite{markovsky2023data}.

\section{Local Data Linearization -- Selective DeePC}
\label{sec:selective_deepc}

The \textbf{Selective DeePC} method \cite{choosewisely} addresses the fundamental limitation of Regularized DeePC when applied to highly nonlinear plants: operating on the full data library forces the algorithm to treat the system as a globally averaged LTI model, introducing significant modeling errors as the robot configuration changes. Its essence is the construction of a local non-parametric model at each control iteration by selecting only those columns of the data matrix that most closely resemble the current system state and the given control objective.

Denoting each column of the Hankel matrix as $\boldsymbol{h}_i$, the database vector and current state vector are defined as:
\begin{itemize}
    \item \textbf{Database vector} ($\boldsymbol{h}_i$) -- the $i$-th column of the Hankel matrix, consisting of historical inputs, outputs, and future outputs:
    \begin{equation}
        \boldsymbol{h}_i = 
        \begin{bmatrix}
        \boldsymbol{u}_{{ini}, i} \\
            \boldsymbol{y}_{{ini}, i} \\
            \boldsymbol{y}_{f, i}
        \end{bmatrix}
    \end{equation}
    
    \item \textbf{Current state vector} ($\boldsymbol{\tilde{h}}$) -- constructed at each control step from current measurements and the reference trajectory:
    \begin{equation}
        \boldsymbol{\tilde{h}} = 
        \begin{bmatrix}
            \boldsymbol{u}_{{ini},~{current}} \\
            \boldsymbol{y}_{{ini},~{current}} \\
            \boldsymbol{r}_{{future}}
        \end{bmatrix}
    \end{equation}
\end{itemize}

At each sampling step, the selection proceeds as follows:
\begin{enumerate}
    \item Calculation of the Euclidean distance $d_i$ between $\boldsymbol{\tilde{h}}$ and each column of the data library:
    \begin{equation}
        d_i = \| \boldsymbol{h}_i - \boldsymbol{\tilde{h}} \|, \quad \text{for } i = 1, \dots, M
    \end{equation}
    \item Sorting indices $i$ by increasing $d_i$.
    \item Selecting $N_{cols}$ nearest-neighbor columns best representing the current dynamics.
    \item Construction of a local Hankel matrix $\mathcal{\tilde H}$ from the selected columns.
\end{enumerate}

The matrix $\mathcal{\tilde H}$ is substituted into (\ref{eq:regularized}) in place of the full data matrix, reducing the dimension of $\boldsymbol{g}$ to $N_{cols}$ and improving control quality through local approximation of the nonlinear dynamics.
    \chapter{Research Results}
\label{ch:rezultaty}

\section{Control System Integration and Plant Definition}
\label{sec:integracja_systemu}

The architecture of the research environment constitutes a synthesis of the elements developed in the previous parts of this work. According to the general closed-loop control scheme (Figure \ref{fig:najwazniejszy_graf}), the system consists of the nonlinear dynamic model of the parallel robot derived in Chapter \ref{ch:model_dynamiczny} and the DeePC controller whose structure is defined in Chapter \ref{ch:deepc}.

\subsection{System Input and Output Definition}

From the perspective of the DeePC regulator, the system is defined solely by the available input and output signals. The control input vector $\boldsymbol{u}$ of dimension $m=2$ corresponds to the torques generated by the actuators in the active joints:
\begin{equation}
    \boldsymbol{u} =
    \begin{bmatrix}
        \tau_1 \\ \tau_4
    \end{bmatrix} \in \mathbb{R}^2
    \label{eq:input_definition}
\end{equation}
The measurement output vector $\boldsymbol{y}$ of dimension $p=2$ contains only the angular positions of the active arms:
\begin{equation}
    \boldsymbol{y} =
    \begin{bmatrix}
        \theta_1 \\ \theta_4
    \end{bmatrix} \in \mathbb{R}^2
    \label{eq:output_definition}
\end{equation}
This deliberately limited observability — excluding velocity measurements and direct end-effector localization — constitutes a more demanding test case, forcing the controller to reconstruct system behavior solely from the input-output history contained in the Hankel matrix.

\section{Control Quality Assessment Methodology}

To quantitatively evaluate control quality, the following indicators were defined \cite{ethzurich, choosewisely}:
\begin{itemize}
    \item \textbf{Mean Absolute Error} ($p$):
    \begin{equation}
        \label{eq:precision}
        p = \frac{1}{N_{sim}} \sum_{k=0}^{N_{sim}} |\boldsymbol{y}_k - \boldsymbol{y}_{ref,k}|
    \end{equation}

    \item \textbf{Maximum Control Error} ($p_{max}$):
    \begin{equation}
        \label{eq:max_precision}
        p_{max} = \max_{k} \left( |\boldsymbol{y}_k - \boldsymbol{y}_{ref,k}| \right)
    \end{equation}

    \item \textbf{Total Control Cost} ($c_{total}$):
    \begin{equation}
        c_{total} = \sum_{k=0}^{N_{sim}} \left( \| \boldsymbol{y}_k - \boldsymbol{y}_{ref,k} \|_{\boldsymbol{Q}}^2 + \| \boldsymbol{u}_k \|_{\boldsymbol{R}}^2 \right)
    \end{equation}

    \item \textbf{Average Optimization Time} ($t_{solve}$) and \textbf{Total Computational Time} ($t_{sel}$) for Selective DeePC:
    \begin{equation}
        t_{sel} = t_{solve} + \frac{1}{N_{sim}} \sum_{k=0}^{N_{sim}} t_{sort,k}
    \end{equation}
\end{itemize}

The indicators $p$ and $p_{max}$ were defined analogously for the task space as $p_{pos}$ and $p_{pos,max}$, computed as the Euclidean distance between the reference and actual end-effector coordinates obtained from forward kinematics.

\section{Research Environment and Implementation}
\label{sec:srodowisko}

All simulations were conducted in MATLAB on a platform equipped with a 12th Gen Intel Core i7-12650H processor and 32 GB of RAM. The YALMIP library \cite{yalmip} and the MOSEK solver \cite{mosek} were used to solve the optimization problem.

The base framework was taken from \cite{zhang_deepc_code, github1}. The most important original contributions include:
\begin{itemize}
    \item \textbf{Implementation of Selective DeePC} using the methodology of \cite{choosewisely}.
    \item \textbf{Mosaic Hankel matrix construction} from multiple independent experiments via (\ref{eq:mosaic_hankel}).
    \item \textbf{Backtrack mechanism} for safe data generation within workspace boundaries.
    \item \textbf{Adaptation to the parallel robot}, including singular configuration avoidance constraints in the QP formulation.
\end{itemize}

\section{Data Preparation and Optimization Formulation}
\label{sec:generacja_danych}

The offline historical database was generated using a random walk method adapted from \cite{choosewisely}. Fifty independent experiments of 100 steps each were conducted, yielding a total dataset of $T=5000$ samples combined into a mosaic Hankel matrix via (\ref{eq:mosaic_hankel}). The exciting signal followed:
\begin{equation}
    \boldsymbol{u}_k = a \cdot \boldsymbol{u}_{k-1} + b \cdot \boldsymbol{n}_k
\end{equation}
with parameters $a=1$, $b=0.1$, $n_a=0.4$, where $\boldsymbol{n}_k$ is drawn from a uniform distribution on $[-n_a, n_a]$. A backtrack mechanism ($n_{back}=10$) ensured all collected trajectories remain kinematically feasible. The resulting point cloud is shown in Figure \ref{fig:A_CHMURA}.

\begin{figure}[!htbp]
    \centering
    \includegraphics[width=0.9\linewidth]{rysunki/4.rezultaty/Chmura_punktów.pdf}
    \caption{Visualization of the training data point cloud against the reference trajectory in the robot workspace}
    \label{fig:A_CHMURA}
\end{figure}

The physical actuator and geometric constraints were defined as:
\begin{align}
    u_{i} &\in [-0.1,\ 0.1]\,\mathrm{Nm}, \quad i=1,2 \\
    y_{1} &\in \left[\tfrac{\pi}{8},\ \tfrac{11}{24}\pi\right], \quad
    y_{2} \in \left[-\tfrac{11}{24}\pi,\ -\tfrac{\pi}{8}\right]
\end{align}
All signals were normalized to $[-1, 1]$ via:
\begin{equation}
    \label{eq:normalization}
    x_{n} = 2\frac{x - x_{min}}{x_{max} - x_{min}} - 1
\end{equation}
Hard box constraints were imposed on both inputs and outputs in the QP task:
\begin{align}
    \boldsymbol{H}_u \boldsymbol{u}_k &\leq \boldsymbol{h}_u, \quad
    \boldsymbol{H}_y \boldsymbol{y}_k \leq \boldsymbol{h}_y
\end{align}
with $\boldsymbol{H}_u = [\boldsymbol{I}; -\boldsymbol{I}]_{4\times2}$ and $\boldsymbol{h}_u = \boldsymbol{1}_{4\times1}$, analogously for outputs. A rate constraint $|\Delta \boldsymbol{u}_k| \leq \Delta u_{max} = 0.4$ was additionally imposed to limit control aggressiveness.

\section{Experimental Configuration}
\label{sec:konfiguracja_eksperymentow}

Controller parameters were tuned on Trajectory A (Lemniscate of Gerono — a smooth $C^\infty$ curve) and subsequently evaluated on Trajectory B. The common simulation parameters were:
\begin{equation}
    T_s = 0.01\,\mathrm{s}, \quad T_{sim}=5\,\mathrm{s}, \quad T_{ini}=8, \quad N=20
\end{equation}
\begin{equation}
    \boldsymbol{Q} = 300 \cdot \boldsymbol{I}_{p}, \quad \boldsymbol{R} = 0.01 \cdot \boldsymbol{I}_{m}
\end{equation}

\textbf{Trajectory B (Polygon)} was the primary evaluation scenario (Figure \ref{fig:robotrownoleglyB}, Table \ref{tab:wspolrzedne_B}). Its vertices are traversed at a constant time of $1\,\mathrm{s}$ per segment ($100$ steps at $T_s=0.01\,\mathrm{s}$), imposing variable velocities and velocity discontinuities at corners. This physically unrealizable reference was intentionally designed to stress-test algorithm robustness under extreme dynamics.

\begin{figure}[!htbp]
    \centering
    \includegraphics[width=1\linewidth]{rysunki/4.rezultaty/druga_trajektoria.pdf}
    \caption{Kinematic scheme of the robot with active link constraints and Trajectory B (Polygon)}
    \label{fig:robotrownoleglyB}
\end{figure}

\begin{table}[!htbp]
    \centering
    \caption{Vertex coordinates of reference Trajectory B}
    \label{tab:wspolrzedne_B}
    \begin{tabular}{ccc}\toprule
        \textbf{Vertex No.} & \textbf{Coordinate} $x$ [m] & \textbf{Coordinate} $y$ [m]\\\midrule
        0 (Start) & $x_0$ & $y_0$ \\
        1 & $0.06$ & $0.00$ \\
        2 & $0.10$ & $0.01$ \\
        3 & $0.20$ & $-0.01$ \\
        4 & $0.16$ & $-0.06$ \\
        5 & $0.06$ & $0.00$ \\\bottomrule
    \end{tabular}
\end{table}

\section{Results -- Trajectory B, No Measurement Noise}

Controller parameters for the noise-free case:
\begin{itemize}
    \item \textbf{Regularized DeePC:} $\lambda_g=100,\ \lambda_{\sigma}=10^8$
    \item \textbf{Selective DeePC:} $\lambda_g=0.1,\ \lambda_{\sigma}=10^8,\ N_{cols}=155$
\end{itemize}

\subsection{Workspace Trajectory Analysis}

The tracking results in the XY plane are presented in Figure \ref{fig:bez_szumu_porownanie_b}. The corner cutting visible at vertices is not an algorithm error but a physically necessary behavior: since the reference requires an instantaneous velocity direction change, the DeePC controller generates a realizable trajectory at the cost of a temporary geometric deviation.

\begin{figure}[!htbp]
    \centering
    \includegraphics[width=0.9\linewidth]{rysunki/4.rezultaty/BEZ SZUMU/Trajektoria_B_bez_szumu.pdf}
    \caption{Trajectory B tracking in the workspace (no noise)}
    \label{fig:bez_szumu_porownanie_b}
\end{figure}

The Selective DeePC method shows a clear advantage in corner precision — the local linearization mechanism enables faster adaptation to direction changes. The Regularized controller exhibits overshoots at vertices, a consequence of the averaging effect introduced by the full historical database. In the second-to-last segment (approx. 3--4 s), both methods show local quality degradation, consistent with a data shortage for this specific motion region near singular configurations.

\subsection{Quantitative Assessment}

\begin{table}[!htbp]
    \centering
    \caption{Joint space control quality indicators -- Trajectory B (no noise)}
    \label{tab:bledy_katy_B}
    \renewcommand{\arraystretch}{1.2}
    \begin{tabular}{c w{c}{2.8cm} w{c}{2.8cm} w{c}{2.8cm} w{c}{2.8cm}}
        \toprule
        \multirow{2}{*}{\textbf{Method}} &
        \multicolumn{2}{c}{\textbf{Tracking Precision $p$} [$\cdot 10^{-3}$ rad]} &
        \multicolumn{2}{c}{\textbf{Maximum Error $p_{max}$} [$\cdot 10^{-3}$ rad]} \\
        \cmidrule(lr){2-3} \cmidrule(lr){4-5}
        & $\theta_1$ & $\theta_4$ & $\theta_1$ & $\theta_4$ \\
        \midrule
        Regularized & $1.96$ & $2.39$ & $8.58$ & $8.20$ \\
        Selective   & $\mathbf{1.77}$ & $\mathbf{1.28}$ & $\mathbf{5.80}$ & $\mathbf{7.07}$ \\
        \bottomrule
    \end{tabular}
\end{table}

\begin{table}[!htbp]
    \centering
    \caption{End-effector positioning error indicators -- Trajectory B (no noise)}
    \label{tab:bledy_pos_B}
    \renewcommand{\arraystretch}{1.2}
    \begin{tabular}{c c c}
        \toprule
        \textbf{Method} &
        \textbf{Mean Error $p_{pos}$} [mm] &
        \textbf{Maximum Error $p_{pos,max}$} [mm] \\
        \midrule
        Regularized & $0.54$ & $2.08$ \\
        Selective   & $\mathbf{0.42}$ & $\mathbf{2.03}$ \\
        \bottomrule
    \end{tabular}
\end{table}

The Selective method dominates across all indicators. The mean angular error for $\theta_4$ is nearly twice as small as in the Regularized method, and the task-space precision is superior in both mean and maximum error (Table \ref{tab:bledy_pos_B}). The Euclidean error profile (Figure \ref{fig:bez_szumu_blad_euklidesowy_b}) confirms that outside the problematic 3--4 s segment, Selective maintains a consistently lower error level.

\begin{figure}[!htbp]
    \centering
    \includegraphics[width=0.9\linewidth]{rysunki/4.rezultaty/BEZ SZUMU/Porownanie_Blad_Euklidesowy_B_bez_szumu.pdf}
    \caption{End-effector Euclidean error over time -- Trajectory B (no noise)}
    \label{fig:bez_szumu_blad_euklidesowy_b}
\end{figure}

\subsection{Control Signal Analysis}

The control signal comparison (Figure \ref{fig:bez_szumu_sygnal_b}) reveals the fundamental operational difference between the two methods. Selective DeePC generates signals with larger amplitude changes, enabling faster cornering response at the cost of chattering. The Regularized method produces smoother profiles, tempered by the global data matrix, at the expense of geometric fidelity.

\begin{figure}[!htbp]
    \centering
    \includegraphics[width=0.9\linewidth]{rysunki/4.rezultaty/BEZ SZUMU/Porownanie_wejscia_B_bez_szumu.pdf}
    \caption{Comparison of control signals for Trajectory B (no noise)}
    \label{fig:bez_szumu_sygnal_b}
\end{figure}

\section{Results -- Trajectory B, With Measurement Noise}
\label{sec:nastawy_szum}

Additive white Gaussian noise was introduced to the joint angle measurements to simulate real encoder imperfections:
\begin{equation}
    \tilde{y}_{i,k} = y_{i,k} + \xi_{i,k}, \quad \xi_{i,k} \sim \mathcal{N}(0, \sigma_i^2)
\end{equation}
with standard deviation $\sigma_i = 0.005 \cdot (y_{max,i} - y_{min,i})$, corresponding to $0.5\%$ of the full operating range. Controller parameters were retuned accordingly:
\begin{itemize}
    \item \textbf{Regularized DeePC:} $\lambda_g=20000,\ \lambda_{\sigma}=10^8$
    \item \textbf{Selective DeePC:} $\lambda_g=200,\ \lambda_{\sigma}=5\cdot10^5,\ N_{cols}=155$
\end{itemize}

\subsection{Workspace Trajectory Analysis}

The noisy trajectory comparison (Figure \ref{fig:z_szumem_porownanie_b}) reveals a fundamental divergence in algorithm behavior:

\begin{figure}[!htbp]
    \centering
    \includegraphics[width=0.9\linewidth]{rysunki/4.rezultaty/Z SZUMEM/Trajektoria_B_szum.pdf}
    \caption{Comparison of polygon shape reproduction in the presence of noise}
    \label{fig:z_szumem_porownanie_b}
\end{figure}

\begin{itemize}
    \item \textbf{Regularized DeePC:} The corner cutting phenomenon is drastically intensified. The robot fails to reach the trajectory turnaround points, tracing curvilinear segments inscribed within the polygon. The bias is significant and variable, resulting from the algorithm averaging noisy data across the full historical database.
    \item \textbf{Selective DeePC:} Despite the noise, this method maintains significantly higher geometric fidelity. Through the first three segments the robot follows the straight lines closely and approaches the vertices. Stability problems appear only in the 4th and 5th segments.
\end{itemize}

\subsection{Quantitative Assessment}

\begin{table}[!htbp]
    \centering
    \caption{Joint space control quality indicators -- Trajectory B (with noise)}
    \label{tab:bledy_katy_B_szum}
    \renewcommand{\arraystretch}{1.2}
    \begin{tabular}{c w{c}{2.8cm} w{c}{2.8cm} w{c}{2.8cm} w{c}{2.8cm}}
        \toprule
        \multirow{2}{*}{\textbf{Method}} &
        \multicolumn{2}{c}{\textbf{Tracking Precision $p$} [$\cdot 10^{-2}$ rad]} &
        \multicolumn{2}{c}{\textbf{Maximum Error $p_{max}$} [$\cdot 10^{-2}$ rad]} \\
        \cmidrule(lr){2-3} \cmidrule(lr){4-5}
        & $\theta_1$ & $\theta_4$ & $\theta_1$ & $\theta_4$ \\
        \midrule
        Regularized & $2.75$ & $2.36$ & $\mathbf{5.59}$ & $7.52$ \\
        Selective   & $\mathbf{2.23}$ & $\mathbf{1.72}$ & $6.95$ & $\mathbf{6.98}$ \\
        \bottomrule
    \end{tabular}
\end{table}

\begin{table}[!htbp]
    \centering
    \caption{End-effector positioning error indicators -- Trajectory B (with noise)}
    \label{tab:bledy_pos_B_szum}
    \renewcommand{\arraystretch}{1.2}
    \begin{tabular}{c c c}
        \toprule
        \textbf{Method} &
        \textbf{Mean Error $p_{pos}$} [mm] &
        \textbf{Maximum Error $p_{pos,max}$} [mm] \\
        \midrule
        Regularized & $\mathbf{6.15}$ & $\mathbf{17.61}$ \\
        Selective   & $6.85$ & $46.04$ \\
        \bottomrule
    \end{tabular}
\end{table}

The mean Cartesian error of Selective is inflated by a local instability in the 4th segment, where oscillations introduce a large initial bias at the start of the 5th segment, visible as the $46.04\,\mathrm{mm}$ peak in Figure \ref{fig:z_szumem_blad_euklidesowy_b}. Outside this critical region, Selective maintains an error below $5\,\mathrm{mm}$ for most of the trajectory, compared to a $10\,\mathrm{mm}$ baseline for Regularized. This highlights the trade-off characteristic of the selective approach: superior mean performance with a risk of local destabilization in data-sparse regions.

\begin{figure}[!htbp]
    \centering
    \includegraphics[width=0.9\linewidth]{rysunki/4.rezultaty/Z SZUMEM/Porownanie_Blad_Euklidesowy_B_szum.pdf}
    \caption{End-effector Euclidean error -- Trajectory B (with noise)}
    \label{fig:z_szumem_blad_euklidesowy_b}
\end{figure}

\subsection{Control Signal Analysis}

The control signal profiles (Figure \ref{fig:z_szumem_sygnal_b}) confirm the contrasting strategies: Selective DeePC performs continuous small corrections to maintain geometric fidelity, while Regularized DeePC generates smooth but geometrically inaccurate outputs. This behavior is consistent across both noise conditions and constitutes the defining operational characteristic of each method.

\begin{figure}[!htbp]
   \centering
   \includegraphics[width=0.9\linewidth]{rysunki/4.rezultaty/Z SZUMEM/Porownanie_wyjsc_B_szum.pdf}
   \caption{Control signal profiles -- Trajectory B (with noise)}
   \label{fig:z_szumem_sygnal_b}
\end{figure}

\section{Computational Performance}

Computation times are summarized in Tables \ref{tab:czasy_bez_szumu} and \ref{tab:czasy_szum}. The results are consistent across both noise conditions, confirming that noise presence does not affect computational complexity.

\begin{table}[!htbp]
    \centering
    \caption{Average computation times -- Trajectory B (no noise)}
    \label{tab:czasy_bez_szumu}
    \renewcommand{\arraystretch}{1.4}
    \begin{tabular}{c c}
        \toprule
        \textbf{Method} & \textbf{Trajectory B} \\
        \midrule
        Regularized DeePC & $t_{solve} = 108.38~\mathrm{ms}$ \\
        Selective DeePC & $t_{sel} = \mathbf{14.65}~\mathrm{ms}$ \quad ($t_{sort} = 1.57~\mathrm{ms}$) \\
        \bottomrule
    \end{tabular}
\end{table}

\begin{table}[!htbp]
    \centering
    \caption{Average computation times -- Trajectory B (with noise)}
    \label{tab:czasy_szum}
    \renewcommand{\arraystretch}{1.4}
    \begin{tabular}{c c}
        \toprule
        \textbf{Method} & \textbf{Trajectory B} \\
        \midrule
        Regularized DeePC & $t_{solve} = 99.96~\mathrm{ms}$ \\
        Selective DeePC & $t_{sel} = \mathbf{13.64}~\mathrm{ms}$ \quad ($t_{sort} = 1.55~\mathrm{ms}$) \\
        \bottomrule
    \end{tabular}
\end{table}

The Selective method is approximately 7 times faster in both conditions. The overhead from the data selection stage ($t_{sort} \approx 1.6\,\mathrm{ms}$) is negligible compared to the gain from reducing the decision variable dimension from $\approx3650$ to $N_{cols}=155$. These results position Selective DeePC as the more promising candidate for real-time implementation, with computation times approaching the $50\,\mathrm{Hz}$ control regime.

\section{Hyperparameter Sensitivity Analysis}
\label{sec:wrazliwosc}

Sensitivity analysis was performed on noisy Trajectory B as the most demanding test case, using the nominal parameter set from Section \ref{sec:nastawy_szum} as the reference.

\subsection{Impact of Regularization Coefficient $\lambda_g$}

\begin{figure}[!htbp]
    \centering
    \includegraphics[width=0.9\linewidth]{rysunki/4.rezultaty/WPŁYW ZMIANY PARAMETRÓW/Analiza_wpływu_lambda_g.pdf}
    \caption{Impact of regularization parameter $\lambda_g$ on control errors}
    \label{fig:wplyw_lambda_g}
\end{figure}

The analysis (Figure \ref{fig:wplyw_lambda_g}) reveals clearly distinct stability thresholds for the two methods:
\begin{itemize}
    \item \textbf{Selective DeePC} requires $\lambda_g > 10^2$ for stable operation. Above $10^3$, a gradual increase in control cost is observed, indicating a well-defined optimal range.
    \item \textbf{Regularized DeePC} requires significantly stronger regularization at $\lambda_g > 10^4$ — a threshold two orders of magnitude higher than Selective.
\end{itemize}
This difference reflects the fundamental distinction between the two approaches: by operating on a local, pre-filtered subset of data, Selective DeePC inherently requires less regularization to suppress noise, as the data selection step already acts as an implicit filter.

\subsection{Impact of Selected Column Number $N_{cols}$ (Selective DeePC)}

\begin{figure}[!htbp]
    \centering
    \includegraphics[width=0.9\linewidth]{rysunki/4.rezultaty/WPŁYW ZMIANY PARAMETRÓW/Analiza_wpływu_Ncols_Selective.pdf}
    \caption{Impact of $N_{cols}$ on Selective DeePC performance}
    \label{fig:wplyw_ncols}
\end{figure}

The analysis (Figure \ref{fig:wplyw_ncols}) reveals three distinct operating regimes:
\begin{itemize}
    \item \textbf{Precision optimum:} The lowest cost function value is recorded at $N_{cols}=170$, where the local model best approximates the robot dynamics from the most relevant historical data.
    \item \textbf{Gradual degradation:} In the range $170$--$275$ columns, a steady increase in control cost is observed as progressively less representative trajectories are included.
    \item \textbf{Instability threshold:} Beyond $N_{cols}=275$, a rapid loss of stability occurs with errors reaching hundreds of millimeters, confirming that including distant, poorly matched data is more harmful than operating on a smaller but well-selected sample.
\end{itemize}
Computation time increases monotonically with $N_{cols}$, consistent with the growing QP problem dimension, reinforcing the importance of selecting the smallest $N_{cols}$ that preserves control quality.

\section{Discussion}

The simulation studies confirm a clear performance differentiation between the two DeePC variants on the nonlinear parallel robot benchmark.

Regularized DeePC exhibits a non-zero systematic bias directly attributable to the global data approach: since the five-bar linkage dynamics are highly nonlinear, the full Hankel matrix cannot simultaneously represent all workspace regions with equal fidelity. The slack variable $\boldsymbol{\sigma}$ partially compensates for this, but at the cost of tracking accuracy. This limitation is significantly amplified under noise, manifesting as the characteristic polygon-to-curve degradation observed in Figure \ref{fig:z_szumem_porownanie_b}.

Selective DeePC addresses this limitation through local dynamics linearization. The improvement in geometric fidelity — even under noise — confirms the hypothesis that restricting the data matrix to nearest-neighbor trajectories reduces the approximation error of the implicit LTI assumption. The cost of this strategy is a higher sensitivity to data-sparse regions: when the local model is poorly conditioned, temporary instability can occur, as observed in the 4th trajectory segment.

A crucial practical finding is the sevenfold reduction in computation time achieved by Selective DeePC, from approximately $100\,\mathrm{ms}$ to $14\,\mathrm{ms}$. This reduction — driven purely by the smaller QP dimension — brings real-time implementation within reach without sacrificing, and in most cases improving, tracking quality.

The sensitivity analysis confirms that parameters tuned on Trajectory A transfer successfully to the more demanding Trajectory B, suggesting that the optimal hyperparameter configuration is primarily governed by robot dynamics rather than specific trajectory excitation. The $\lambda_g$ threshold analysis further underscores the structural advantage of the selective approach: its implicit noise filtering through data selection relaxes the regularization requirements by two orders of magnitude relative to the global method.
    \chapter{Conclusions}
\label{ch:podsumowanie}

This paper presented the implementation and comparative evaluation of Regularized and Selective DeePC applied to position control of a nonlinear five-bar parallel robot. The cut-joint method proved an effective and modular framework for deriving joint-space dynamics of closed kinematic chains. The following conclusions are drawn from the simulation studies:

\begin{enumerate}
    \item \textbf{Data quality dependency:} Under ideal conditions, both controllers achieve high tracking precision. The introduction of measurement noise causes significant quality degradation in both variants, confirming that data-driven methods are critically sensitive to the quality of the offline Hankel matrix dataset.

    \item \textbf{Regularized DeePC:} This method operates conservatively, generating smooth control signals with low rate of change. While beneficial for real hardware implementation, this filtering effect introduces steady-state bias and corner cutting on sharp trajectory turns, a limitation that is substantially amplified under noise.

    \item \textbf{Selective DeePC:} Local linearization via nearest-neighbor data selection reproduces sharp trajectory changes with significantly higher geometric fidelity and lower mean tracking error. The trade-off is reduced stability in data-sparse workspace regions, manifesting as control signal oscillations and occasional local destabilization.

    \item \textbf{Computational performance:} The reduction of the QP problem dimension to $N_{cols}$ yields a sevenfold decrease in computation time relative to the Regularized method (from $\approx 100\,\mathrm{ms}$ to $\approx 14\,\mathrm{ms}$), positioning Selective DeePC as the more viable candidate for real-time control applications.

    \item \textbf{Hyperparameter transferability:} Parameters tuned on a smooth training trajectory successfully generalize to the more demanding polygon test case, suggesting that the optimal configuration is primarily governed by robot dynamics rather than specific trajectory excitation.
\end{enumerate}

\section{Future Work}

The results motivate the following directions for further research:

\begin{enumerate}
    \item \textbf{Experimental verification:} Deployment of Selective DeePC on a physical laboratory setup, requiring real training data collection and validation of noise robustness under hardware disturbances.
    \item \textbf{Stabilization of the selective method:} Development of continuity conditions for the data selection algorithm to prevent abrupt changes in the local Hankel matrix structure, with the aim of eliminating the observed oscillatory behavior in data-sparse regions.
    \item \textbf{Implementation optimization:} Evaluation of alternative QP solvers (e.g., OSQP) and lower-level code translation to verify strict real-time feasibility.
\end{enumerate}

    % Bibliografia - musi być, jest w pliku EE-dyplom.bib
    % Bibliography - must exist, is in the file EE-dyplom.bib
    \bibliografia

    % Strony końcowe - można zakomentować, jeśli zbędne
    % Additional pages - comment out if not needed
    
    % Wykaz symboli i skrótów - patrz opis w tekście przykładowym
    %\glsaddall
    %\acronymslist
    % Spis rysunków
    \listoffigures
    % Spis tabel
    \listoftables
    % Załączniki (plik appendices.tex)
    \easyappendices
\end{document}
%%%%%%%%%%%%%%%%%%%%%%%%%%%%%%%%%%%%%%%%%%%%%%%%%%%%%%%%%%%%%%%%%%%%%%%%%%%

